%% LyX 2.4.4 created this file.  For more info, see https://www.lyx.org/.
%% Do not edit unless you really know what you are doing.
\documentclass[english,t]{beamer}
\usepackage{mathptmx}
\usepackage[T1]{fontenc}
\usepackage[latin9]{inputenc}
\usepackage{amssymb}
\usepackage{cancel}

\makeatletter
%%%%%%%%%%%%%%%%%%%%%%%%%%%%%% Textclass specific LaTeX commands.
% this default might be overridden by plain title style
\newcommand\makebeamertitle{\frame{\maketitle}}%
% (ERT) argument for the TOC
\AtBeginDocument{%
  \let\origtableofcontents=\tableofcontents
  \def\tableofcontents{\@ifnextchar[{\origtableofcontents}{\gobbletableofcontents}}
  \def\gobbletableofcontents#1{\origtableofcontents}
}

%%%%%%%%%%%%%%%%%%%%%%%%%%%%%% User specified LaTeX commands.
\usetheme{Warsaw}
% or ...

\setbeamercovered{transparent}
% or whatever (possibly just delete it)
\setbeamercovered{invisible} %makes overlays invisible


%gets rid of navigation symbols
\setbeamertemplate{navigation symbols}{}

\usepackage{multicol}

\usepackage{tikz}
\usetikzlibrary{plotmarks}
\usepackage{pgfplots}
\pgfplotsset{compat=newest}

\usepackage{calc}
\usepackage[nomessages]{fp}

\usepackage{advdate}    % Advancing/saving dates
\usepackage{datetime}   % Dates formatting
\usepackage{datenumber} % Counters for dates

%\usepackage{pgfpages}
%\pgfpagesuselayout{4 on 1}[a4paper, landscape, border shrink=7mm]
%\pgfpageslogicalpageoptions{1}{border code=\pgfusepath{stroke}}
%\pgfpageslogicalpageoptions{2}{border code=\pgfusepath{stroke}}
%\pgfpageslogicalpageoptions{3}{border code=\pgfusepath{stroke}}
%\pgfpageslogicalpageoptions{4}{border code=\pgfusepath{stroke}}
%%%Don't forget to add 'handout'

\makeatother

\usepackage{babel}
\begin{document}
\newcommand{\xmax}{2000}
\newcommand{\xmin}{0}
\newcommand{\xstep}{0} %must be xmin+n
\newcommand{\ymax}{500}
\newcommand{\ymin}{0}
\newcommand{\ystep}{0} %must be ymin+n

\newcounter{countEx} \setcounter{countEx}{0}
\newcounter{countProb} \setcounter{countProb}{0}
\resetcounteronoverlays{countEx} \resetcounteronoverlays{countProb}

\newcommand{\ex}[3]{
	\addtocounter{countEx}{1}
	\only<#1-#2>{\begin{exampleblock} {Example \chNum.\secNum.\arabic{countEx}.} #3	\end{exampleblock}}
}

\newcommand{\prob}[4]{
	\addtocounter{countProb}{1}
	\only<#1-#2>{\begin{exampleblock} {Problem  \chNum.\secNum.\arabic{countProb}.\,#3} #4	\end{exampleblock}}
}

\beamerdefaultoverlayspecification{<+->}

%%%%Chapter and Section numbers%%%%%
\newcommand{\chNum}{1}
\newcommand{\secNum}{2}
\newcommand{\secTitle}{Linear Functons and Applications}

\title[Section \chNum.\secNum: \secTitle]{Section \chNum.\secNum: \secTitle}

\subtitle{Finite Mathematics~~MTH 132~~Spring 2014}

\author{Dr.\,James Ashe}

\titlegraphic{\includegraphics[height=1.5cm]{/home/jim/JamesAshe/JCSU/images/jcsu_bull}}

\institute{Johnson C. Smith University}

\date{{}}

\makebeamertitle 
\begin{frame}{Linear Functions and Applications}

\begin{columns}


\column{5.75cm}

Linear functions are just linear equations.
\begin{itemize}
\item <2->Equations are functions.
\item <3->Function notation is just substitution.
\end{itemize}
\ex{4}{}{Let $g(x)=-4x+2$. Find $g(2),g(-1),g(0),g(x+1)$.}

\column{6cm}

\end{columns}
\end{frame}

\begin{frame}{Word Problems (Yay!)}

\begin{block}{General Word Problem Strategies}

\begin{enumerate}
\item <2->Read the problem carefully. Identify a question (usually at the
end).
\item <3->Represent an \emph{unknown }with an $x$. 

\begin{enumerate}
\item <4->Write all other unknowns in terms of $x$.
\end{enumerate}
\item <5->Write an equation(s) relating $x$ to the \emph{known} values.
\item <6->Solve the equation.
\item <7->Use your solution to answer \emph{all }the questions.
\item <8->Check!
\end{enumerate}
\end{block}
\end{frame}

\begin{frame}{Cost Function}

\ex{1}{3}{The marginal cost of using a firm charges \$10 \textbf{per hour} plus a \textbf{flat rate} of \$50. Find the cost function, $C(x)$, of using this moving firm for $x$ hours.}

\ex{4}{}{A moving firm charges \$10 \textbf{per hour } plus an additional \textbf{service charge} of \$50. Find the cost function, $C(x)$, of using this moving firm for $x$ hours.}
\begin{itemize}
\item <2->$C(x)=mx+b$, so find $m$ and $b$.
\item <3->$m$ is the \emph{marginal cost}, and \emph{$b$ }is the \emph{fixed
cost}.
\item <4->\emph{maginal cost=cost per $x$=each $x$ costs}
\end{itemize}
\prob{5}{}{Find the cost function and identify all variables used.}{The internet site "Music.com" charges a \$10 registration fee plus 99 cents per song.}
\end{frame}

\begin{frame}{Revenue and Profit Function}

A small publishing company has fixed costs of \$525, and can print
a book for \$2. The same book sells for \$5 a piece.
\begin{columns}


\column{5.5cm}

\ex{1}{1}{Find the cost function.}

\ex{2}{2}{Find the revenue function.}

\ex{3}{3}{Find the profit funciton.}

\ex{4}{4}{How many books must be produced and sold in order to break even?}

\ex{5}{5}{How many books must be produced and sold in order to make a profit of \$1,000?}
\begin{block}<6->{Profit Funtion}
$P(x)=R(x)-C(x)$
\end{block}

\column{6cm}

\end{columns}
\end{frame}

\begin{frame}{Revenue and Profit Function}

A small publishing company has fixed costs of \$525, and can print
a book for \$2. The same book sells for \$5 a piece.

\renewcommand{\xmax}{600}
\renewcommand{\xmin}{0}
\renewcommand{\xstep}{0} %must be xmin+n
\renewcommand{\ymax}{1200}
\renewcommand{\ymin}{-600}
\renewcommand{\ystep}{0} %must be ymin+n
%%%%%%%
\begin{tikzpicture} 
\begin{axis}[height=5.5cm, width=5.5cm, axis lines=middle,minor x tick num={4},minor y tick num={4},grid=both,%xtick={0,50,...,250},%ytick={0,10,20,...,40},%axis line style={-latex}, 
xmin=\xmin,xmax=\xmax, ymin=\ymin,ymax=\ymax,
%restrict y to domain=-1:10,%no markers, 
	xlabel={$q$},ylabel={$p$}, 
	%every axis x label/.append style={anchor=north}, 
	%every axis y label/.append style={anchor=east}
	]
  
	\only<1-2>{\addplot[domain=\xmin:\xmax,thick,draw=red,->,color=red]{5*x} node[pos=.45,left] {$R(x)$};}
	\only<2-2>{\addplot[domain=\xmin:\xmax,thick,draw=blue,color=blue,->]{525+2*x} node[pos=.75,above] {$C(x)$};}
	\only<3->{\addplot[domain=\xmin:\xmax,thick,draw=blue,->]{5*x-525-2*x} node[pos=.7,above] {$P(x)$};}
	%\only<9->{\addplot[color=blue,solid,mark=*,style={font=\footnotesize}] coordinates {(1,-1)} node[right] {$(1,-1)$};}
	%\only<10->{\addplot[color=blue,solid,mark=*,style={font=\footnotesize}] coordinates {(2,1)} node[above,right] {$(2,1)$};}
	
\end{axis}
\end{tikzpicture}
\end{frame}

\begin{frame}{Supply and Demand}

\ex{1}{4}{Suppose that the demand and price for strawberries are related by \[ p=D(q)=5-0.25q \] Find the price at the level of demand of 0 and 400 quarts}

\ex{5}{6}{Suppose that the supply and price for strawberries are related by \[ p=S(q)=0.25q \] Find the price at the level of supply of 0 and 400 quarts}
\begin{itemize}
\item <2-6>$p$ is the price (in dollars).
\item <3-4>$q$ is the quantity \emph{demanded} (in hundereds of quarts).
\item <4-4>$D(q)$ is price at which consumers will buy $q$ number of
products.
\item <5-6>$q$ is the quantity \emph{supplied }(in hundreds of quarts).
\item <6-6>$S(q)$ is the price at which suppliers will produce $q$ number
of products.
\end{itemize}
\end{frame}

\begin{frame}{Supply and Demand}

\ex{1}{}{Find the \emph{equilibrum quantity} for the supply, $S(q)=0.25q$, and demand, $D(q)=5-0.25q$ functions.}
\begin{columns}


\column{5.75cm}

\vspace{-.5cm}
%%%%%%%
\begin{tikzpicture} 
\begin{axis}[height=5.5cm, width=5.5cm, axis lines=middle,minor x tick num={4},minor y tick num={4},grid=both,%xtick={0,50,...,250},%ytick={0,10,20,...,40},%axis line style={-latex}, 
xmin=\xmin,xmax=\xmax, ymin=\ymin,ymax=\ymax,
%restrict y to domain=-1:10,%no markers, 
	xlabel={$q$},ylabel={$p$}, 
	%every axis x label/.append style={anchor=north}, 
	%every axis y label/.append style={anchor=east}
	]
  
	\only<2->{\addplot[domain=\xmin:\xmax,thick,draw=red,->]{.25*x} node[pos=.65,above] {$S(q)$};}
	\only<3->{\addplot[domain=\xmin:\xmax,thick,draw=blue,->]{500-.25*x} node[pos=.65,above] {$D(q)$};}
	%\only<9->{\addplot[color=blue,solid,mark=*,style={font=\footnotesize}] coordinates {(1,-1)} node[right] {$(1,-1)$};}
	%\only<10->{\addplot[color=blue,solid,mark=*,style={font=\footnotesize}] coordinates {(2,1)} node[above,right] {$(2,1)$};}
	
\end{axis}
\end{tikzpicture}

\column{5.5cm}

\end{columns}
\end{frame}

\begin{frame}{Supply and Demand}

\begin{columns}


\column{5.5cm}

\ex{1}{1}{Let the supply and demand functions for a gourmet choclate bar be $p=S(q)=\frac{5}{4}q$ and $p=D(q)=70-\frac{3}{4}q$, where $p$ is the price in dollars and $q$ is the number of bars. Find the equilibrium price and quantity.}

\prob{2}{2}{}{Let the supply and demand functions for sugar be $p=S(q)=1.4q-0.6$ and $p=D(q)=-2q+3.2$, where $p$ is the price per pound in dollars and $q$ is quantity in thousands of pounds. Graph both functions on the same axes; find the equilibrium price and quantity.}

\column{6cm}

\end{columns}
\end{frame}

\begin{frame}{Summary}

\begin{block}{}

\begin{enumerate}
\item Cost function, $C(x)$.

\begin{enumerate}
\item Find the marginal and fixed cost.
\item Find the cost function.
\end{enumerate}
\item Profit function, $P(x)=R(x)-C(x)$.

\begin{enumerate}
\item Find the profit function.
\item Find the break even point.
\item Find quantity needed to make a desired profit.
\end{enumerate}
\item Supply and Demand, $S(q)$ and $D(q)$.

\begin{enumerate}
\item Find the price for a given quantity of demand/supply.
\item Find the equilibrium quantity.
\item Find the equilibrium price.
\end{enumerate}
\end{enumerate}
\end{block}
\end{frame}

\end{document}
